\section{Related Work}
\label{sec:related}

\subsection{Digital Policing and Inter-Agency Sharing}

CCTNS and ICJS are the baseline systems for the Indian context. Official
sources state that CCTNS is in use across all 17,798 police stations and that
ICJS links the police/CCTNS pillar with courts, prisons, forensics, and
prosecution~\cite{pib_cctns_2026,mha_icjs_2026}. The public NCRB Crime in
India data is useful for understanding registered crime at an aggregate level,
but it is not an individual-level access-log or FIR dataset for this paper's
problem~\cite{ncrb_2023}. India-focused AI/crime studies use such public data
for spatial or statistical analysis, but they do not model classified
inter-agency access governance~\cite{sharma_india_crime_2023,naick_murder_xai_2024,reza_crix_2025}.

Recent legal research also warns that AI use in criminal justice raises data
protection questions in India, especially because sensitive criminal-justice
data can involve accused persons, suspects, victims, witnesses, and the wider
public~\cite{anbarasi_ai_criminal_justice_2026}. This supports the paper's
scope: the work does not replace CCTNS or ICJS, and it does not treat police
records as ordinary open data. It asks what kind of secure overlay can sit
above a CCTNS/ICJS-compatible environment when sensitive records must be requested,
explained, and reviewed across agencies.

\subsection{Blockchain Audit and Evidence Management}

Blockchain has already been studied for provenance, access control, and digital
evidence management. NISTIR~8403 explains why blockchain can support
tamper-resistant access-control records, while also noting practical concerns
such as governance, scalability, and system design choices~\cite{nist_blockchain_access_2022}.
Hyperledger Fabric is a useful reference for a permissioned setting where the
participants are known organizations rather than anonymous public
miners~\cite{androulaki_fabric_2018}. Certificate Transparency also provides a
non-blockchain reference point for append-only Merkle-tree audit logs and
consistency proofs~\cite{laurie_ct_2013}. This is important because the paper
does not assume that a blockchain is automatically better than every signed log.

Digital evidence systems and lawful-evidence chains show how blockchains can
support chain-of-custody, provenance, off-chain evidence storage, and forensic
integrity~\cite{kim_two_level_2021,li_lechain_2021,jeong_fabric_evidence_2020}.
Other work studies blockchain-based auditable access control, privacy-preserving
access control, forensic log access, and digital evidence protection
~\cite{akhtar_auditable_access_2020,rouhani_access_control_survey_2019,hu_accountable_private_ac_2024,islam_forensic_log_2023,banerjee_dashdikpala_2026}.
Blockchain has also been used in smart public-safety systems where distributed
participants need traceability and auditability~\cite{xu_blendsps_2020}.

These works are relevant, but they do not fully solve the problem studied in
this paper. Evidence provenance is not the same as deciding whether a sensitive
inter-agency access request should be allowed, denied, or escalated. A ledger
can show that a record of a decision was not changed, but it cannot by itself
prove that the decision was correct, privacy-preserving, or properly explained.

\subsection{Access Control for Sensitive Records}

RBAC is a mature baseline for organizing permissions around roles
~\cite{sandhu_rbac_1996}. However, role-only access is too coarse for the
setting studied here. ABAC allows authorization to depend on subject, object,
action, and environmental attributes~\cite{nist_abac_2014}. XACML is also
important as a policy-language reference because it formalizes policy decisions,
obligations, advice, and policy evaluation structures~\cite{oasis_xacml_2013}.
For police data sharing, access may depend on jurisdiction, purpose, record
sensitivity, approval status, credential validity, case assignment, and
emergency context.

Fabric-based ABAC work shows that blockchain infrastructure and access control
can be combined with off-chain encrypted storage patterns
~\cite{zhao_fabric_abac_2022}. Privacy-preserving ABAC work shows that even
attribute disclosure during authorization can create privacy concerns, motivating
selective-disclosure, homomorphic-encryption, and zero-knowledge approaches
~\cite{kerl_privacy_abac_2025}. The present paper does not implement a full
cryptographic privacy-preserving ABAC protocol. Instead, it treats metadata
exposure as an evaluation issue and keeps raw sensitive records off-chain.

However, generic ABAC or blockchain-ABAC work usually does not include a
police-specific allow/deny/escalate workflow. It also does not usually treat
the explanation itself as an artifact that must be logged and checked later.
This is the gap used in this paper: access control is necessary, but it must be
connected to audit evidence and explanation evidence.

\subsection{Explainable AI and High-Stakes Review}

The XAI part of this paper is not added for presentation. It is part of the
audit model. Rudin argues that high-stakes decisions should prefer
interpretable models where this is feasible~\cite{rudin_2019}. Recent work on
law-enforcement XAI also highlights that explanations must support different
stakeholders and must not encourage blind trust in automated systems
~\cite{zocholl_xai_law_enforcement_2025}. Work on human-centered explanation
shows that explanations are social and contrastive, not just feature lists
~\cite{miller_xai_social_2019}. Counterfactual explanations are also relevant
because an access requester may need to know what missing condition would have
changed a denial or escalation decision~\cite{wachter_counterfactual_2018}.

Standard XAI methods such as LIME and SHAP explain model predictions through
local surrogate explanations and feature attributions
~\cite{ribeiro_lime_2016,lundberg_shap_2017}. This paper does not depend on a
black-box predictor for the main access decision; it uses a deterministic policy
oracle so that labels are transparent in the first benchmark. Still, the XAI
literature is useful for defining explanation artifacts, explanation coverage,
counterfactual validity, and evaluation discipline~\cite{doshi_velez_interpretability_2017}.

There is also work that directly connects blockchain and XAI. Nassar et al.
propose blockchain support for trustworthy and explainable AI systems
~\cite{nassar_blockchain_xai_2020}; Malhotra et al. study blockchain audit
trailing of XAI decisions using IPFS and Ethereum~\cite{malhotra_xai_audit_2021};
and legal/XAI studies discuss blockchain-based audit of legal or justice-related
AI decisions~\cite{sachan_legal_ai_blockchain_2024,demertzis_xai_justice_2023}.
These works support the idea of logging explanation evidence, but they do not
provide the CCTNS/ICJS-compatible police access-governance benchmark studied here.

\subsection{Why This Is Not Predictive Policing}

This paper avoids individual crime prediction and suspect prediction. Prior
work has shown that predictive-policing systems can create feedback loops when
future data is influenced by earlier policing activity~\cite{ensign_feedback_2018}.
Algorithmic policing work also emphasizes assumptions, evaluation, and
accountability problems in prediction-led policing~\cite{bennett_moses_policing_2018}.
Machine-learning crime prediction has a large literature, but the datasets,
labels, and evaluation conditions vary widely~\cite{jenga_crime_prediction_2023}.
Criminal-justice prediction work also shows that complex models are not
automatically better than simpler baselines~\cite{dressel_fairness_limits_2018}.
Fairness literature shows that some fairness criteria cannot all be satisfied
simultaneously under common conditions~\cite{chouldechova_2017,kleinberg_tradeoffs_2016}.

These findings support a cautious framing. The paper uses AI to explain access
decisions and detect policy drift in logs, not to predict criminals or crime
locations.

\subsection{Research Gap}

The existing literature gives useful building blocks: blockchain audit,
permissioned ledgers, ABAC, off-chain storage, and XAI for high-stakes
settings. The missing piece is their joint evaluation in an inter-agency access
governance workflow for CCTNS/ICJS-compatible police record requests. In particular,
the literature does not sufficiently study what happens when a log is validly
re-signed after policy output is corrupted. In that case, ordinary integrity
checks may pass even though the decision is wrong. This paper studies that gap
with a synthetic benchmark and compares defenses under clear visibility
assumptions.
